%% sec_platform.tex
%% ─────────────────────────────────────────────────────────────────────────────
%% Section 3 — Target Platform: AMD Versal AI Engine Architecture
%% ─────────────────────────────────────────────────────────────────────────────

\section{Target Platform: AMD Versal AI Engine Architecture}\label{sec:platform}

\subsection{Device overview}

The AMD Versal AI Core Series VCK190 evaluation kit is built around the
XCVC1902 Adaptive Compute Acceleration Platform (ACAP).  The device integrates
three compute domains on a single die: a dual-core ARM Cortex-A72 application
processor (1.7\,GHz) used for system control and host orchestration; a
dual-core ARM Cortex-R5F real-time processor; and the programmable logic (PL)
fabric with 899\,840 LUTs, 1\,968 DSP58 engines, and 158\,MB of embedded
memory (UltraRAM + BRAM).  The AI Engine array is the key resource for
computation in this work.

\subsection{AI Engine tile microarchitecture}\label{sec:aie_tile}

The device contains \textbf{400 \aie tiles} arranged in a $50\times 8$ array
and clocked at 1\,GHz.  Each tile comprises:
\begin{itemize}
  \item A 512-bit SIMD vector unit supporting single-precision (float32) and
    fixed-point arithmetic;
  \item A 32-bit scalar RISC core for control flow;
  \item 32\,KB of local data memory (plus 128\,KB addressable via neighbours);
  \item \textbf{16\,KB of dedicated program memory (PM)}---the hard constraint
    that drives the architecture of this work.
\end{itemize}

\subsection{Cascade interface}\label{sec:cascade}

Adjacent tiles in the same row are connected by a \textbf{cascade interface}:
a 384-bit unidirectional accumulator bus clocked in lock-step with the vector
pipeline.  Key properties relevant to this work are:
\begin{itemize}
  \item \textbf{Deterministic, zero-overhead data transfer}: one 384-bit word
    ($4\times\texttt{cfloat}$) per clock cycle, no flow-control signals.
  \item \textbf{Hardware FIFO}: a 4-entry FIFO decouples producer and consumer
    by up to 4 cycles, absorbing minor skew without stalling.
  \item \textbf{Fixed direction per row}: even rows (0, 2, 4, 6) cascade
    left-to-right; odd rows (1, 3, 5, 7) cascade right-to-left.  Kernel
    placement must match this hardware topology.
  \item \textbf{Bandwidth}: $384\,\text{bits}/\text{cycle} \times 1\,\text{GHz}
    = 48\,\text{GB/s}$, far exceeding the stream I/O bandwidth from the shim.
\end{itemize}

\subsection{Stream switching and PLIO}

Tiles communicate with the PL through \emph{Programmable Logic I/O} (PLIO)
ports at the shim row.  On XCVC1902, valid shim columns for PLIO placement are
columns 6--43 (columns 0--5 and above 43 host AIE-PL interface tiles but not
direct PLIO sites).  The stream switch network routes packets between tiles
using packet headers, enabling \texttt{pktsplit} and \texttt{pktmerge}
primitives that distribute or collect one stream to/from multiple tiles without
dedicated PL logic.

\subsection{Program memory as the primary constraint}

Unlike GPU or CPU codebases where code size is virtually unconstrained, the
16\KB PM limit is hard and non-negotiable.  A kernel that overflows PM cannot
be compiled.  Table~\ref{tab:pm_footprint} quantifies the PM contribution of
each HELAS function used in \ggttg, establishing why no single tile can hold
the complete computation.

\begin{table}[h]
\caption{Estimated program memory footprint of each HELAS function used for
\ggttg on the \aie.  Values are representative; exact figures depend on
compiler optimisation flags.  The 16\KB hardware limit is shown for
reference.}
\label{tab:pm_footprint}
\begin{tabular}{@{}lrl@{}}
\toprule
HELAS function / component & PM (bytes) & Used in \\
\midrule
\texttt{VVV1P0\_1} (off-shell gluon)          & ${\sim}2{,}000$ & K1 \\
\texttt{FFV1P0\_3} (off-shell boson)          & ${\sim}2{,}700$ & K3 (moved) \\
\texttt{FFV1\_0/1/2} (FFV amplitudes)         & ${\sim}4{,}400$ & K2a, K2b \\
\texttt{VVV1\_0} (triple-gluon amplitude)     & ${\sim}2{,}100$ & K3 \\
\texttt{VVVV1/3/4P0\_1} (4-gluon contact)    & ${\sim}6{,}600$ & K4 \\
Kernel logic + token serialisation            & ${\sim}11{,}000$ & all \\
Softfloat runtime                             & ${\sim}1{,}500$ & all \\
\midrule
\textbf{Monolithic total}                     & ${\sim}38{,}000$ & --- \\
\textbf{16\KB hardware limit}                 & $16{,}384$ & --- \\
\bottomrule
\end{tabular}
\end{table}

The \textbf{VVV Exclusion Principle} emerges directly from
Table~\ref{tab:pm_footprint}: \texttt{VVV1P0\_1} (${\sim}2.0\KB$) and
\texttt{VVV1\_0} (${\sim}2.1\KB$), when combined with the kernel logic
overhead, cannot coexist in a single 16\KB tile.  They must reside in
separate kernels---\texttt{K1} and \texttt{K3} respectively.
